\documentclass[11pt,a4paper]{article}

\newcommand{\docTitle}{Rejestr Ryzyk}

% ─────────────────────────────────────────────────────────────────
% Preambuła wspólna dla dokumentów projektu Feedback
% ─────────────────────────────────────────────────────────────────

% Kodowanie i język
\usepackage[utf8]{inputenc}
\usepackage[T1]{fontenc}
\usepackage[polish]{babel}

% Układ strony
\usepackage{geometry}
\geometry{margin=2.5cm}

% Tabele
\usepackage{booktabs}
\usepackage{tabularray}
\usepackage{longtable}
\usepackage{array}

% Grafika i strony specjalne
\usepackage{afterpage}
\usepackage{rotating}
\usepackage{pdflscape}

% Kolory
\usepackage{xcolor}
\definecolor{accent}{RGB}{30, 60, 100}
\definecolor{lowrisk}{RGB}{200, 230, 200}
\definecolor{medrisk}{RGB}{255, 235, 170}
\definecolor{highrisk}{RGB}{245, 195, 195}

% Tytuły, listy, akapity
\usepackage{titlesec}
\usepackage{parskip}
\usepackage{enumitem}
\usepackage{fancyhdr}

% Hyperref i cleveref – kolejność ma znaczenie:
% hyperref musi być załadowany przed cleveref
\usepackage{hyperref}
\hypersetup{colorlinks=true, linkcolor=accent, urlcolor=accent, citecolor=accent}
\usepackage[nameinlink]{cleveref}
% Polskie nazwy dla cleveref – formy skrócone unikają problemu z odmianą
\crefname{section}{sek.}{sek.}
\Crefname{section}{Sek.}{Sek.}
\crefname{subsection}{sek.}{sek.}
\Crefname{subsection}{Sek.}{Sek.}
\crefname{table}{tab.}{tab.}
\Crefname{table}{Tab.}{Tab.}
\crefname{figure}{rys.}{rys.}
\Crefname{figure}{Rys.}{Rys.}
\crefname{equation}{rów.}{rów.}
\Crefname{equation}{Rów.}{Rów.}
\crefname{enumi}{pkt}{pkt}
\Crefname{enumi}{Pkt}{Pkt}


% Numerowane sekcje z kolorem akcentu
\titleformat{\section}{\large\bfseries\color{accent}}{\thesection.}{0.6em}{}[\titlerule]
\titleformat{\subsection}{\normalsize\bfseries}{\thesubsection}{0.6em}{}
\titleformat{\subsubsection}{\normalsize\itshape}{\thesubsubsection}{0.6em}{}

% Łamanie wierszy – tolerancja dla polskich tekstów
\setlength{\emergencystretch}{3em}
\hyphenpenalty=1000
\tolerance=2000

% ─────────────────────────────────────────────────────────────────
% Nagłówki i stopki – wymaga zdefiniowania \docTitle przed \input
% ─────────────────────────────────────────────────────────────────

\pagestyle{fancy}
\fancyhf{}
\rhead{\small \docTitle}
\lhead{\small Projekt: \textit{Feedback}}
\cfoot{\thepage}
\renewcommand{\headrulewidth}{0.4pt}


% Dodatkowe pakiety/ustawienia specyficzne dla rejestru
\usepackage{makecell}

% Kolorowanie komórek P/W/E – używamy \SetCell{bg=...} bezpośrednio
% (skróty \newcommand nie działają z tabularray, bo \SetCell musi być
% pierwszym tokenem komórki przed ekspansją makra)

\begin{document}
	
	% ── STRONA TYTUŁOWA ──────────────────────────────────────────────
	\begin{titlepage}
		\centering
		\vspace*{2cm}
		{\LARGE\bfseries\color{accent} \docTitle}\\[0.5cm]
		{\large Projekt: \textit{Feedback}}\\[0.3cm]
		{\normalsize Stan na: \today}\\[2cm]
		
		\begin{tabular}{ll}
			\textbf{Kierownik projektu:}   & Kacper Trzciński \\[0.3em]
			\textbf{Wersja dokumentu:}     & 1.0 \\[0.3em]
			\textbf{Źródło operacyjne:}    & Google Sheets (folder \texttt{/Dokumenty}) \\[0.3em]
		\end{tabular}
		
		\vfill
		{\small Zarządzanie Przedsięwzięciami Informatycznymi\\
			Wydział Matematyki i Nauk Informacyjnych, Politechnika Warszawska}
	\end{titlepage}
	
	\pagenumbering{roman}
	
	% ── NOTA WPROWADZAJĄCA ──────────────────────────────────────────
	\section*{O dokumencie}
	\addcontentsline{toc}{section}{O dokumencie}
	
	Niniejszy dokument stanowi formalny rejestr ryzyk projektu \textit{Feedback} w stanie
	na dzień podany na stronie tytułowej. Pełne zasady zarządzania ryzykiem (skale, macierz,
	tolerancje, procedury) opisuje dokument \textit{Plan Zarządzania Ryzykiem}.
	
	Operacyjnym kanałem aktualizacji rejestru jest arkusz Google Sheets w folderze
	\texttt{/Dokumenty} na dysku zespołu --- to tam członkowie zespołu rejestrują zmiany na
	bieżąco. Niniejszy dokument \LaTeX{} jest \textbf{wersjonowanym mirrorem} arkusza,
	odświeżanym przed każdym przeglądem końca sprintu i przy każdej eskalacji do Komitetu
	Sterującego. Historia zmian poszczególnych pozycji prowadzona jest w arkuszu (osobna
	karta) i nie jest powielana w niniejszym dokumencie.
	
	\subsection*{Legenda skrótów}
	\begin{itemize}[noitemsep]
		\item \textbf{Typ}: Z --- zagrożenie, S --- szansa.
		\item \textbf{P} --- prawdopodobieństwo (1--5).
		\item \textbf{W} --- wpływ (1--5).
		\item \textbf{E} --- ekspozycja ($E = P \cdot W$); kolor komórki odpowiada klasie ryzyka (zielony / żółty / czerwony).
		\item \textbf{Status}: \textit{open} --- aktywne, \textit{mitigated} --- po wdrożeniu reakcji, \textit{materialised} --- zmaterializowane, \textit{closed} --- zamknięte.
		\item \textbf{Rezydualne} --- ocena P/W/E po wdrożeniu zaplanowanej reakcji; ,,---'' gdy reakcja nie została jeszcze wdrożona lub nie zmieniła oceny.
	\end{itemize}
	
	\subsection*{Podsumowanie}
	\begin{itemize}[noitemsep]
		\item Łączna liczba pozycji: 18 (16 zagrożeń, 2 szanse).
		\item Ryzyka wysokie (czerwone): 2.
		\item Ryzyka średnie (żółte): 11.
		\item Ryzyka niskie (zielone): 5.
	\end{itemize}
	
	\clearpage
	\pagenumbering{arabic}

% ── REJESTR RYZYK (CUSTOM ROZMIAR - JEDNA DŁUGA STRONA) ─────────
% Szerokość = 29.7cm (A4 landscape), wysokość = 80cm.
% Pakiet geometry w \newgeometry nie obsługuje paperwidth/paperheight,
% więc używamy obejścia: ustawiamy wymiary papieru PRZED \newgeometry i
% jednocześnie przekazujemy total widoczne textwidth/textheight,
% a do PDF wpisujemy fizyczny rozmiar przez \pdfpage*.
\pdfpagewidth=60cm
\pdfpageheight=40cm
\paperwidth=60cm
\paperheight=40cm
\newgeometry{total={56cm,35.7cm}, top=2cm, left=2cm, includehead=false, includefoot=false}

	% ── AKTUALIZACJA SZEROKOŚCI NAGŁÓWKA DLA FANCYHDR ──
% Przypisujemy aktualną szerokość tekstu (\textwidth) do nagłówka
\setlength{\headwidth}{\textwidth}

\clearpage
\setlength{\tabcolsep}{6pt} % Zwiększono nieco odstęp dla czytelności na tak dużej stronie

\begin{landscape}
	\section*{Rejestr ryzyk}
	\addcontentsline{toc}{section}{Rejestr ryzyk}
	
	\begin{longtblr}[
		caption = {Rejestr ryzyk projektu \textit{Feedback}.},
		label = {tab:rejestr},
		]{
			width = \linewidth, % Tabela rozciąga się na pełną szerokość zdefiniowanej strony (116cm)
			colspec = {
				|p{1cm}       % ID
				|p{2.2cm}     % Data (lekki zapas na format rrrr-mm-dd)
				|c            % Typ (krótki znak Z/S)
				|X[4, l]      % Opis (bardzo szeroki, wyrównany do lewej)
				|X[1.8, l]    % Właściciel
				|c|c|c        % P, W, E
				|X[1.8, l]    % Reakcja
				|X[4, l]      % Akcje minimalizujące / zwiększające
				|p{2cm}       % Termin
				|X[1.5, l]    % Status
				|p{2.2cm}|    % Rezydualne
			},
			rowhead = 1,
			row{1}  = {font=\bfseries\small, c, m},
			rows    = {font=\small, halign=l, valign=t},
			hlines, vlines,
		}
		ID & Data & \makecell{Typ} & Opis (zdarzenie / przyczyna / skutek) & Właściciel & P & W & E & Reakcja & Akcje minimalizujące / zwiększające & Termin & Status & \makecell{Rezy-\\dualne} \\
		% ────── RYZYKA TECHNICZNE ──────
		R-01 & 2026-04-15 & Z &
		Mechanika parowania okazuje się trudna do zbalansowania (sterowanie nieintuicyjne, brak ,,zabawowości''). Przyczyna: złożoność fizyki parowania, niewielkie doświadczenie zespołu w game-feel. Skutek: konieczność przeprojektowania kluczowej mechaniki w późnej fazie projektu. &
		A. Płocharski &
		4 & \SetCell{bg=highrisk} 5 & \SetCell{bg=highrisk} 20 &
		Ograniczanie &
		Wczesny grywalny prototyp parowania (sprint 2). Cykliczne playtesty wewnętrzne co sprint. Playtest zewnętrzny po prototypie ,,szarej kostki''. Zarezerwowany budżet czasu na iteracje mechaniki w sprintach 3--5. &
		Sprint 2 & open & 3 / 4 / 12 \\
		
		R-02 & 2026-04-15 & Z &
		Wybrany silnik (np. Godot) okazuje się niewystarczający dla wymagań mechaniki parowania --- problemy z wydajnością fizyki lub brakujące funkcje. Przyczyna: niedostateczna walidacja silnika na etapie wyboru. Skutek: migracja na inny silnik lub gruntowna przebudowa rozwiązań. &
		F. Shramko &
		2 & \SetCell{bg=highrisk} 5 & \SetCell{bg=medrisk} 10 &
		Ograniczanie &
		Spike techniczny w sprincie 1: implementacja minimalnego prototypu parowania w wybranym silniku przed pełnym commitmentem. Identyfikacja punktów ryzyka technicznego i alternatyw (Unity, custom). &
		Sprint 1 & open & 1 / 5 / 5 \\
		
		R-03 & 2026-04-15 & Z &
		Niska wydajność (spadki FPS poniżej zakładanego minimum) na docelowym sprzęcie. Przyczyna: złożoność symulacji fizyki parowania, nieoptymalne shadery. Skutek: pogorszenie jakości oddania, możliwa redukcja zakresu wizualnego. &
		F. Shramko &
		3 & \SetCell{bg=medrisk} 3 & \SetCell{bg=medrisk} 9 &
		Ograniczanie &
		Profilowanie wydajności w każdym sprincie implementacyjnym. Budżet FPS na pojedyncze efekty. Plan B: redukcja gęstości cząstek pary i jakości shaderów do akceptowalnego poziomu (degradacja \textit{Could}). &
		Sprint 4 & open & --- \\
		
		R-04 & 2026-04-15 & Z &
		Bugi krytyczne wykryte po finalnej kompilacji, niemożliwe do naprawienia przed prezentacją. Przyczyna: brak czasu na pełne testy regresyjne. Skutek: prezentacja gry z widocznymi defektami. &
		K. Trzciński &
		3 & \SetCell{bg=medrisk} 4 & \SetCell{bg=medrisk} 12 &
		Ograniczanie &
		Code freeze na 2 tygodnie przed prezentacją końcową. Dedykowane sprinty testowe (sprint 7 i 8). Rotacja testerów wewnątrz zespołu. Lista bugów blokujących utrzymywana w systemie zadań z etykietą \texttt{blocker}. &
		Sprint 7 & open & 2 / 4 / 8 \\
		
		R-05 & 2026-04-15 & Z &
		Utrata danych projektu (kod, assety, dokumentacja). Przyczyna: błąd lokalny, awaria dysku, błąd w obsłudze repozytorium. Skutek: utrata postępu, opóźnienie. &
		R. Głasek &
		2 & \SetCell{bg=medrisk} 4 & \SetCell{bg=lowrisk} 8 &
		Ograniczanie &
		Kod w repozytorium GitHub (gałąź główna chroniona, wymagane PR). Assety i dokumentacja w Google Drive (wersjonowanie automatyczne). Cotygodniowy lokalny backup wykonywany przez właściciela ryzyka. &
		Sprint 1 & mitigated & 1 / 3 / 3 \\
		
		% ────── RYZYKA ZESPOŁOWE ──────
		R-06 & 2026-04-15 & Z &
		Choroba lub niedostępność członka zespołu w krytycznym momencie (egzaminy, praca, sytuacje losowe). Przyczyna: równoległe obciążenie studiami i pracą zawodową. Skutek: opóźnienie zadań, presja na pozostałych. &
		K. Trzciński &
		4 & \SetCell{bg=medrisk} 3 & \SetCell{bg=medrisk} 12 &
		Ograniczanie &
		Wymóg sparowania (peer review) każdego zadania --- każdy obszar zna co najmniej 2 osoby. Bufor planistyczny: zespół planuje sprint na 80\% pojemności. Wczesne komunikowanie planowanej niedostępności (urlopów, sesji). &
		Sprint 1 & open & 3 / 3 / 9 \\
		
		R-07 & 2026-04-15 & Z &
		Konflikt w zespole prowadzący do spadku produktywności lub odejścia członka zespołu. Przyczyna: różnice w zaangażowaniu, niejasny podział odpowiedzialności. Skutek: opóźnienia, atmosfera utrudniająca pracę. &
		K. Trzciński &
		2 & \SetCell{bg=medrisk} 4 & \SetCell{bg=lowrisk} 8 &
		Ograniczanie &
		Retrospektywa po każdym sprincie z przestrzenią na zgłaszanie problemów. Jasna definicja ról i odpowiedzialności w karcie projektu. Otwarty kanał z kierownikiem projektu dla zgłoszeń niepublicznych. &
		Ciągły & open & --- \\
		
		R-08 & 2026-04-15 & Z &
		Brak kompetencji w obszarze krytycznym (np. shadery, optymalizacja, projektowanie poziomów). Przyczyna: zespół studencki, brak doświadczenia w gamedev. Skutek: niska jakość komponentu lub konieczność rezygnacji. &
		A. Płocharski &
		3 & \SetCell{bg=medrisk} 3 & \SetCell{bg=medrisk} 9 &
		Ograniczanie &
		Identyfikacja luk kompetencyjnych na początku semestru. Czas na naukę w sprincie 1 i 2. Konsultacje z prowadzącym i społecznością (np. Godot Discord). Plan B: redukcja zakresu obszaru. &
		Sprint 2 & open & 2 / 3 / 6 \\
		
		R-09 & 2026-04-15 & Z &
		Przeciążenie kierownika projektu (kumulacja zadań koordynacyjnych i implementacyjnych). Przyczyna: brak delegacji, koncentracja odpowiedzialności. Skutek: pogorszenie jakości koordynacji, opóźnienia w dokumentacji. &
		K. Trzciński &
		3 & \SetCell{bg=medrisk} 3 & \SetCell{bg=medrisk} 9 &
		Ograniczanie &
		Wyznaczony zastępca KP (rotacyjnie). Delegacja prowadzenia stand-upów. Wyraźny budżet czasu KP na koordynację (max 30\% jego czasu w sprincie). &
		Sprint 2 & open & --- \\
		
		% ────── RYZYKA PROCESOWE ──────
		R-10 & 2026-04-15 & Z &
		Niedoszacowanie zadań w planowaniu sprintu prowadzące do regularnego niedowożenia. Przyczyna: brak doświadczenia w estymacji, optymistyczne planowanie. Skutek: kumulacja opóźnień, narastający dług. &
		K. Trzciński &
		4 & \SetCell{bg=medrisk} 3 & \SetCell{bg=medrisk} 12 &
		Ograniczanie &
		Estymacja w story points (porównawcza, nie godzinowa). Po sprincie 2 kalibracja velocity zespołu. Bufor 20\% w planowaniu. Retrospektywa estymacji co sprint. &
		Sprint 2 & open & 3 / 3 / 9 \\
		
		R-11 & 2026-04-15 & Z &
		Dokumentacja pisana w ostatniej chwili, niespójna ze stanem projektu. Przyczyna: priorytet implementacji, brak nawyku dokumentowania na bieżąco. Skutek: dokumenty oddane w słabej jakości, obniżenie oceny. &
		A. Bielicki &
		4 & \SetCell{bg=medrisk} 3 & \SetCell{bg=medrisk} 12 &
		Ograniczanie &
		Definicja ukończenia (DoD) zadania obejmuje aktualizację dokumentacji. Sekcja ,,dokumentacja'' w przeglądzie końca sprintu. Wspólny szablon \LaTeX{} (preamble.tex) zmniejszający narzut. &
		Sprint 1 & open & 3 / 3 / 9 \\
		
		R-12 & 2026-04-15 & Z &
		Niejasne wymagania od prowadzącego, zmieniające się w trakcie projektu. Przyczyna: ewolucja oczekiwań w trakcie semestru, interpretacje. Skutek: przebudowy zakresu, marnowanie pracy. &
		K. Trzciński &
		3 & \SetCell{bg=medrisk} 3 & \SetCell{bg=medrisk} 9 &
		Ograniczanie &
		Pytania do prowadzącego konsolidowane raz na sprint (jeden kanał, jeden moderator). Decyzje prowadzącego dokumentowane w rejestrze decyzji projektu. Wyjaśnienia w razie wątpliwości na bieżąco. &
		Ciągły & open & --- \\
		
		R-13 & 2026-04-15 & Z &
		Niezaliczenie etapu projektu (Karta projektu, prezentacja koncepcji, WBS + harmonogram). Przyczyna: niespełnienie wymagań prowadzącego, niska jakość dokumentu. Skutek: utrata punktów, ryzyko niezaliczenia przedmiotu. &
		K. Trzciński &
		2 & \SetCell{bg=highrisk} 5 & \SetCell{bg=medrisk} 10 &
		Ograniczanie &
		Konsultacja draftu z prowadzącym przed oddaniem (minimum 1 tydzień przed terminem). Wewnętrzne review dokumentu przez 2 członków zespołu spoza autora. Lista checkpointów na podstawie wymagań przedmiotu. &
		Każdy etap & open & 1 / 5 / 5 \\
		
		R-14 & 2026-04-15 & Z &
		Brak playtesterów zewnętrznych w kluczowych momentach (early prototype, beta). Przyczyna: brak wcześniejszych ustaleń, niska dostępność wolontariuszy. Skutek: brak zewnętrznej walidacji mechaniki, ryzyko rozjazdu z odbiorcą. &
		T. Herman &
		3 & \SetCell{bg=medrisk} 3 & \SetCell{bg=medrisk} 9 &
		Ograniczanie &
		Zbudowanie listy 8--10 potencjalnych playtesterów (znajomi z roku, niższe roczniki) w sprincie 1. Kalendarz sesji playtestowych z wyprzedzeniem. Plan B: playtest wewnątrzzespołowy z rotacją ról. &
		Sprint 2 & open & 2 / 3 / 6 \\
		
		% ────── RYZYKA ZEWNĘTRZNE ──────
		R-15 & 2026-04-15 & Z &
		Awaria sprzętu członka zespołu w późnej fazie projektu (laptop deweloperski). Przyczyna: awaria losowa, brak rezerwy sprzętowej. Skutek: utrata produktywności dotkniętej osoby na czas naprawy. &
		K. Trzciński &
		2 & \SetCell{bg=medrisk} 3 & \SetCell{bg=lowrisk} 6 &
		Akceptacja &
		Środowisko deweloperskie odtwarzalne z repozytorium (kontener, instrukcja setup). Możliwość pracy na pożyczonym sprzęcie z dostępu MiNI. Backup dzienny do chmury. &
		Sprint 1 & mitigated & 1 / 3 / 3 \\
		
		R-16 & 2026-04-15 & Z &
		Zmiana prowadzącego przedmiot lub jego oczekiwań w trakcie roku akademickiego. Przyczyna: czynniki zewnętrzne (urlop, choroba, decyzje wydziału). Skutek: konieczność dostosowania zakresu lub formy oddania. &
		K. Trzciński &
		1 & \SetCell{bg=medrisk} 3 & \SetCell{bg=lowrisk} 3 &
		Akceptacja &
		Dokumentacja prowadzona zgodnie ze standardami metodyk projektowych (PRINCE2-lite), nie tylko pod konkretnego prowadzącego. Wczesne zorientowanie się w polityce wydziału w razie sygnałów. &
		Ciągły & open & --- \\
		
		% ────── SZANSE ──────
		R-17 & 2026-04-15 & S &
		Mechanika parowania okazuje się na tyle ciekawa, że przyciąga uwagę szerszej społeczności (np. itch.io, Reddit r/gamedev). Przyczyna: unikalność konceptu, jakość prototypu. Skutek: feedback od zewnętrznych odbiorców, potencjalne dalsze rozwijanie po projekcie. &
		A. Płocharski &
		2 & \SetCell{bg=medrisk} 3 & \SetCell{bg=lowrisk} 6 &
		Wykorzystanie &
		Wcześniejsze przygotowanie krótkiego devloga / GIF-a z mechaniki. Publikacja po sprincie 4 (po stabilnym prototypie). Monitorowanie reakcji i włączenie cennego feedbacku do iteracji. &
		Sprint 4 & open & --- \\
		
		R-18 & 2026-04-15 & S &
		Któryś z członków zespołu odkrywa elegancką, niespodziewaną interakcję mechaniki parowania (\textit{emergent gameplay}) podczas implementacji. Przyczyna: właściwości fizycznego systemu, eksperymenty zespołu. Skutek: dodatkowy mechanizm zwiększający głębię rozgrywki bez większych kosztów. &
		T. Herman &
		3 & \SetCell{bg=medrisk} 3 & \SetCell{bg=medrisk} 9 &
		Wykorzystanie &
		Spotkania ,,zabaw z prototypem'' co 2 sprinty --- cały zespół eksperymentuje z mechaniką bez konkretnego zadania. Rejestrowanie odkryć w kanale Discord \texttt{\#emergent}. Włączanie najciekawszych do projektu poziomów. &
		Sprint 3 & open & --- \\
	\end{longtblr}
	\end{landscape}
	
	\pdfpagewidth=210mm
	\pdfpageheight=297mm
	\paperwidth=210mm
	\paperheight=297mm
	\restoregeometry
	
	% ── PRZYWRÓCENIE SZEROKOŚCI NAGŁÓWKA DLA STRON A4 ──
	\setlength{\headwidth}{\textwidth}
	
\end{document}